\centerline{\bf \Huge Теорема Виета}


\textbf{Формулы Виета для квадратного уравнения.}
Если $x_1$ и $x_2$ - корни квадратного уравнения $a x^2+b x+c=0$, то

$$
\left\{\begin{array}{l}
x_1+x_2=-\frac{b}{a}, \\
x_1 x_2=\frac{c}{a} .
\end{array}\right.
$$


\textbf{Формулы Виета для приведённого квадратного уравнения.}
В частном случае, если $a=1$ (приведённая форма $x^2+p x+q=0$ ), то

$$
\left\{\begin{array}{l}
x_1+x_2=-p, \\
x_1 x_2=q .
\end{array}\right.
$$


\textbf{Формулы Виета для кубического уравнения.}
Если $x_1, x_2, x_3$ - корни кубического уравнения $p(x)=a x^3+b x^2+c x+d=0$, то

$$
\left\{\begin{array}{l}
x_1+x_2+x_3=-\frac{b}{a} \\
x_1 x_2+x_1 x_3+x_2 x_3=\frac{c}{a} \\
x_1 x_2 x_3=-\frac{d}{a}
\end{array}\right.
$$

\textbf{Формулы Виета в общем виде.}
Если $x_1, x_2, \ldots, x_n$ - корни многочлена

$$f(x)=x^n+a_1 x^{n-1}+a_2 x^{n-2}+\ldots+a_n$$

в виде симметрических многочленов от корней, а именно:

$$
\left\{\begin{array}{l}
a_1=-\left(x_1+x_2+\ldots+x_n\right) \\
a_2 =x_1 x_2+x_1 x_3+\ldots+x_1 x_n+x_2 x_3+\ldots+x_{n-1} x_n \\
a_3 =-\left(x_1 x_2 x_3+x_1 x_2 x_4+\ldots+x_{n-2} x_{n-1} x_n\right)\\
\qquad\qquad\qquad\qquad\qquad \ldots\\
a_{n-1} =(-1)^{n-1}\left(x_1 x_2 \ldots x_{n-1}+x_1 x_2 \ldots x_{n-2} x_n+\ldots+x_2 x_3 \ldots x_n\right), \\
a_n =(-1)^n x_1 x_2 \ldots x_n .
\end{array}\right.
$$

Иначе говоря, $(-1)^k a_k$ равно сумме всех возможных произведений из $k$ корней. Если многочлен не приведён, что справа делим на старший коэффициент.


\newpage

\hspace{3mm}

\problem{1}
Пусть $a, b, c, d$ - такие вещественные числа, что

$$
a^3+b^3+c^3+d^3=a+b+c+d=0
$$

Докажите, что сумма каких-то двух из этих чисел равна нулю.

\textit{Подсказка: Введите обозначение $a+b=u$. Тогда $c+d=-u$. Преобразуйте сумму кубов.}

\problem{2}
Целые числа $a, b$ и $c$ таковы, что числа $\dfrac{a}{b}+\dfrac{b}{c} +\dfrac{c}{a}$ и $\dfrac{a}{c}+\dfrac{c}{b}+\dfrac{b}{a}$ тоже целые. Докажите, что $|a|=|b|=|c|$.

\textit{Подсказка: введите такие обозначения $x=\dfrac{a}{b},\  y=\dfrac{b}{c}$ и $z=\dfrac{c}{a}$.}


\problem{3$^{*}$}
Натуральные числа $a, b, c, d, e$ и $f$ таковы, что число $S=a+b+c+d+e+f$ делит числа $a b c+d e f$ и $a b+b c+c a-d e-e f-f d$. Докажите, что число $S$ составное.

\textit{Подсказка Рассмотрите следующий многочлен: $f(t)=(t-a)(t-b)(t-c)-(t+d)(t+e)(t+f)$}